\section{Introduction}

Los autómatas celulares son sistemas dinámicos discretos que se han utilizado para modelar sistemas complejos~\cite{Torres2025_complejo,sole1996_cs_caos,Hoekstra2010_cs_ca}. A partir de reglas locales muy simples producen comportamientos que van desde la convergencia a un estado homogéneo, fácilmente predecible, hasta dinámicas caóticas~\cite{wolfram1982_ca,Wolfram2002_nks}. Múltiples investigaciones han propuesto clasificaciones para distinguir estos comportamientos, como las cuatro clases de Wolfram~\cite{Wolfram2002_nks}, la clasificación de Li y Packard~\cite{li1990_structure} y otras más revisadas por Martínez~\cite{genaro2013_class}, o el conjunto de reglas \emph{life-like}~\cite{lifewiki_life}. Estas clasificaciones suelen diseñarse según el comportamiento de las reglas, pero ¿qué tan parecidas son las reglas de una misma clase? y ¿qué tan diferentes son de las reglas de otra?

Actualmente, las clasificaciones son la única medida que podemos utilizar para decidir si dos reglas son parecidas, además de comparar directamente su comportamiento exacto. Sin embargo, ambas dan una respuesta binaria y no existe una medida que nos permita cuantificar qué tan similares son. El propósito de esta investigación es definir un método que permita comparar el comportamiento de cualesquiera dos reglas, aunque tengan distinto número de estados o de vecinos, y medir su grado de similitud.

Para conocer las similitudes, primero identificamos las relaciones exactas que ya existen entre reglas de dominios diferentes (Sección~\ref{sec:inheritance}) y las usamos como referencia para validar los métodos propuestos. Después comparamos la estructura del grafo de transición mediante su forma canónica (Sección~\ref{sec:structural}) y proponemos métodos espectrales que representan cada regla con un vector (Sección~\ref{sec:spectral}): los eigenvalores, los momentos espectrales y los momentos espectrales con masa de atractor. Por último, evaluamos este último método con reglas de dominios distintos, con espacios de distinto tamaño y agrupando las reglas elementales (Sección~\ref{sec:experiments}).
